\section{Calculus}

\subsection{Convergence Tests}

\begin{theorem}[Divergence Test]
The series $\sum_{n \geq 0} a_n$ diverges if $\lim_{n\rightarrow\infty} a_n \neq 0$.
\end{theorem}

\begin{theorem}[$p$-series Test]
The series $\sum_{n \geq 0} \frac{1}{n^p}$, for some $p\in\R$ converges if and only if
$p > 1$.
\end{theorem}

\begin{theorem}
Suppose that $\sum a_n, \sum b_n$ are nonnegative series. Consider the quantity
\[
\lim_{n\to\infty}\frac{a_n}{b_n}=L
\]
where $L\in[0,\infty]$. Then we have the following three conclusions:
\begin{itemize}
    \item If $L=0$, the series $\sum a_n$ converges if $\sum b_n$ converges.
    \item If $L=\infty$, the series $\sum b_n$ converges if $\sum a_n$ converges.
    \item If $L\in(0,\infty)$, the series $\sum a_n$ converges if and only if $\sum b_n$ converges.
\end{itemize}
\end{theorem}

\begin{theorem}
Suppose that $\{a_n\}$ is a positive series such that there exists a continuous,
decreasing function $f(x)$ satisfying $f(n)=a_n$. Then the series $\{a_n\}$ converges 
if and only if the improper integral
\[
\int_0^\infty f(x) dx
\]
exists (converges).
\end{theorem}

\begin{theorem}[Alternating Series Test]
Suppose that $\{a_n\}$ is a positive, (eventually) decreasing series. 
If $\lim_{n\to\infty}a_n=0$, the series $\sum_{n \geq 0}(-1)^n a_n$ converges.
\end{theorem}

\begin{example}
The series $\sum \frac{1}{n}$ doesn't converge, but the alternating 
series $\sum \frac{(-1)^n}{n}$ does converge.
\end{example}

\subsection{Integration, FTC, and Variants}
Recall $u$-substitution. Suppose we have the integral
\[
\int_a^b f(g(x))g'(x)dx.
\]
Then, if we set $u = g(x)$, then this integral becomes
\[
\int_{g(a)}^{g(b)} f(u) du.
\]
Here are some useful specific cases:
\begin{align*}
\int_0^a f\left(\frac{x}{a}\right)dx &= a\int_0^1 f(u)du \\
\int_a^b f(x-c) dx &= \int_{a-c}^{b-c}f(u) du \\
\int_a^b f(a+b-x) dx &= \int_a^b f(u) du
\end{align*}

\begin{theorem}[Leibniz's Rule]
Suppose that $a(x)$ and $b(x)$ are differentiable, and that
$f(x,t)$ and $\frac{\partial f}{\partial x}(x,t)$ are continuous.
Then
\[
\frac{d}{dx}\int_{a(x)}^{b(x)}f(x,t) dt
=
f(x,b(x))b'(x) - f(x,a(x))a'(x) +
\int_{a(x)}^{b(x)}
\frac{\partial f}{\partial x}(x,t) dt.
\]
\end{theorem}


\subsection{Fubini}

The assumptions for Fubini are quite technical, so I leave it up 
to the wikipedia page. This is a result from measure theory which essentially
states when we can swap integrals. For absolutely integrable functions (or Lebesgue 
measurable), we can do this.

\begin{exercise}
Assume $b>a>0$. Calculate the integral:
\[
\int_0^\infty \frac{e^{-ax}-e^{-bx}}{x}\,dx.
\]
\textit{Hint: Write the integrand as an integral from $a$ to $b$.}
\begin{enumerate}[label=(\Alph*)]
    \item $\ln a-\ln b^2$
    \item $\ln b-\ln a^2$
    \item $\ln a^2-\ln b^2$
    \item $\ln a-\ln b$
    \item $\ln b-\ln a$
\end{enumerate}
\end{exercise}
\begin{solution}
(B). The integral is
\[
\int_0^{\infty} \int_{a}^b e^{-xy} dy dx.
\]
We use Fubini to flip the order of integration.
\end{solution}

\subsection{Things to watch out for}

We include some common misconseptions.

\begin{remark}
Let $a \in \R$, and suppose we have some $f(x)$ defined on the reals 
such that $\int_a^{\infty} f(x) dx$ converges. This does not imply that 
$f(x) \rightarrow 0$ as $x \rightarrow \infty$. We can consider some 
`spike function', whose spikes become narrow and narrow as $x$ becomes large.
For instance, let $f(x)$ be zero everywhere, except a spike of hegiht $n$ 
and width $\frac{2}{n^3}$, around every positive integer. Then, the integral 
is $\frac{\pi^2}{6}$, which converges, but clearly the limit doesn't exist. 
If we want something like this, then it must be like the following statements:
\begin{itemize}
\item If the improper integral exists and $\lim_{x \rightarrow \infty} f(x) = L$, 
where $L$ is some finite real, then $L = 0$.
\item If the improper integral exists and $f(x)$ is uniformly continuous, then 
the limit exist (Barbalat's lemma). By the above, it must be zero.
\item If $f$ is nonnegative and eventually nonincreasing, then $\lim_{x \rightarrow \infty} f(x) = 0$.
\end{itemize}
\end{remark}

\begin{remark}
We note that $|f(x)| \rightarrow \infty$  does not imply that $f(x) \rightarrow \infty$
or $f(x)\rightarrow -\infty$.
\end{remark}

\subsection{Standard Calculus Facts to Quicken Things}

We start withs ome basic limits.
\begin{itemize}
\item $\lim_{x\rightarrow 0}\frac{\sin x}{x}=1$
\item $\lim_{x\rightarrow 0}\frac{\tan x}{x}=1$
\item $\lim_{x\rightarrow 0}\frac{1-\cos x}{x^2}=\frac{1}{2}$
\item $\lim_{x\rightarrow 0}\frac{\arctan x}{x}=1$
\end{itemize}

Useful Taylor expansions:
\begin{align*}
\sin x&=x-\frac{x^3}{3!} + \frac{x^5}{5!} - \dots \\
\cos x&=1-\frac{x^2}{2!} + \frac{x^4}{4!} - \dots \\
\tan x&=x+\frac{x^3}{3}+ \frac{2x^5}{15} + \dots.
\end{align*}

\subsection{Group Theory and Binary Operations}

\subsubsection*{Basic Group Facts}

For a finite group $G$,
\[
\operatorname{ord}(g)\mid |G|
\qquad(g\in G).
\]

\textbf{Cauchy's Theorem.}
If a prime $p$ divides $|G|$, then $G$
contains an element of order $p$.

For a cyclic group $G=\langle g\rangle$
of order $n$,
\[
\boxed{
\operatorname{ord}(g^k)
=
\frac{n}{\gcd(n,k)}.
}
\]

\subsubsection*{Permutation Cycles}

The inverse of a cycle reverses its direction:
\[
\boxed{
(a\ b\ c)^{-1}=(a\ c\ b).
}
\]

Disjoint cycles commute.

The order of a permutation is the least
common multiple of its disjoint cycle lengths.

\subsubsection*{Two Binary Operations}

Suppose two binary operations $\cdot$ and
$\star$ share a two-sided identity $e$ and
satisfy the interchange law
\[
(a\cdot b)\star(c\cdot d)
=
(a\star c)\cdot(b\star d).
\]

Then the Eckmann--Hilton argument gives
\[
\boxed{
a\star b=a\cdot b
}
\]
and the common operation is commutative:
\[
\boxed{
a\cdot b=b\cdot a.
}
\]

\textbf{GRE strategy:}
When two operations are related by an
unfamiliar identity, try substituting
identity elements into different positions.


\subsection{Topology: Connectedness and Closure}

\subsubsection*{Basic Implications}

\[
\boxed{
\text{Path-connected}
\implies
\text{Connected}.
}
\]

The converse is false.

A continuous image of a connected set
is connected.

A continuous image of a path-connected set
is path-connected.

If $C$ is connected and
\[
C\subseteq A\subseteq\overline C,
\]
then $A$ is connected.

In particular, the closure of a connected
set is connected.

\subsubsection*{Topologist's Sine Curve}

Consider
\[
S=
\left\{
\left(x,\sin\frac1x\right):0<x\le1
\right\}
\cup
\bigl(\{0\}\times[-1,1]\bigr).
\]

This set is:
\begin{itemize}
    \item connected;
    \item closed;
    \item not path-connected.
\end{itemize}

The infinitely rapid oscillations prevent
a continuous path from connecting the
vertical segment to the oscillating graph.

\textbf{GRE strategy:}
Do not confuse adjoining limit points
with creating continuous paths.


\subsection{Piecewise Smoothness}

Suppose
\[
f(x)=
\begin{cases}
g(x),&x<a,\\
h(x),&x\ge a.
\end{cases}
\]

For continuity at $a$, require
\[
\boxed{g(a)=h(a)}
\]
using the appropriate one-sided extensions.

For $f$ to be $C^1$ across $a$, require
\[
\boxed{
g(a)=h(a),\qquad
g'(a)=h'(a).
}
\]

For $f$ to be $C^2$ across $a$, require
\[
\boxed{
\begin{aligned}
g(a)&=h(a),\\
g'(a)&=h'(a),\\
g''(a)&=h''(a).
\end{aligned}
}
\]

Assume the pieces extend smoothly
to the joining point.

\textbf{GRE strategy:}
Check continuity first. If the function
values do not match, there is no reason
to calculate the derivatives.


\subsection{Geometry: Useful Formulas}

\subsubsection*{Circular Sectors}

For a circle of radius $r$ and angle
$\theta$ measured in radians,
\[
\boxed{A=\frac12r^2\theta.}
\]

The arc length is
\[
\boxed{s=r\theta.}
\]

\subsubsection*{Hyperbolas}

For
\[
\frac{x^2}{a^2}-\frac{y^2}{b^2}=1,
\]
the foci are $(\pm c,0)$, where
\[
\boxed{c^2=a^2+b^2.}
\]

The asymptotes are
\[
\boxed{y=\pm\frac ba x.}
\]

The first-quadrant asymptote makes angle
\[
\boxed{\theta=\arctan\left(\frac ba\right)}
\]
with the positive $x$-axis.

\subsubsection*{Ellipses}

For
\[
\frac{x^2}{a^2}+\frac{y^2}{b^2}=1,
\qquad a>b>0,
\]
the focal distance satisfies
\[
\boxed{c^2=a^2-b^2.}
\]

\subsubsection*{Other Geometric Formulas}

A tetrahedron with three mutually
perpendicular edges of lengths $a,b,c$
has volume
\[
\boxed{V=\frac{abc}{6}.}
\]

For a sphere of radius $r$,
\[
\boxed{
V=\frac43\pi r^3,
\qquad
S=4\pi r^2.
}
\]
